% ============================================================
\chapter{Integer Programming}
\label{ch:integer}
% ============================================================

\section{Introduction}

In linear programming, optimal solutions live in a continuous world: one can assign $3.7$ trucks to a route. But in the real world, a truck is an integer. You cannot cut an airplane in half to serve two routes. This integrality constraint, seemingly innocuous, radically transforms the nature of the problem: integer programming is NP-hard in general, and the solution methods are fundamentally different. The \emph{branch and bound} method, developed by Ailsa Land and Alison Doig in 1960, intelligently explores a tree of subproblems; \emph{Gomory cuts} (1958) add constraints to tighten the continuous relaxation; \emph{branch and cut} combines both. These algorithms are at the heart of modern solvers (CPLEX, Gurobi) that solve industrial problems with millions of variables.

\begin{intuition}
Continuous linear programming may produce fractional solutions (e.g., build 3.7
tables). Naively rounding is generally neither optimal nor feasible. Integer
programming provides systematic methods to find the best integer solution.
\end{intuition}

\section{Formulation}

\begin{definition}[Integer linear program]
\index{integer programming}
An \textbf{integer linear program} (ILP) is:
\[
  \min\; c^\top x \quad \text{s.t.} \quad Ax \leq b, \quad x \in \Z^n_+.
\]
If only some variables are integer, it is a \textbf{mixed-integer program}
(MIP). If all variables are binary ($x \in \{0,1\}^n$), it is a
\textbf{binary program}.
\end{definition}

\begin{definition}[LP relaxation]
\index{LP relaxation}
The \textbf{LP relaxation} of an ILP is the linear program obtained by
replacing $x \in \Z^n_+$ with $x \in \R^n_+$. Its optimal value provides a
\textbf{lower bound} (for minimization) on the ILP optimum.
\end{definition}

\begin{proposition}[Relaxation bound]
Let $z^*_{\mathrm{LP}}$ and $z^*_{\mathrm{IP}}$ be the LP and IP optima.
Then $z^*_{\mathrm{LP}} \leq z^*_{\mathrm{IP}}$. If the LP optimum is
integral, it is also optimal for the ILP.
\end{proposition}

\begin{keyformulas}[title=\textbf{Common binary modeling tricks}]
\begin{itemize}
  \item \textbf{Selection}: $x_j = 1$ if item $j$ is selected.
  \item \textbf{Implication}: $x_i = 1 \Rightarrow x_j = 1$ modeled as
        $x_i \leq x_j$.
  \item \textbf{Disjunction}: at least one of $x_1, \ldots, x_k$:
        $x_1 + \cdots + x_k \geq 1$.
  \item \textbf{Big-M}: $y \leq M x$ forces $y = 0$ when $x = 0$.
\end{itemize}
\end{keyformulas}

\section{Branch and Bound}

\begin{definition}[Branch and Bound]
\index{branch and bound}
\textbf{Branch and Bound} (B\&B) solves an ILP by implicit enumeration over
a search tree. At each node, the LP relaxation is solved and pruning rules
eliminate suboptimal branches.
\end{definition}

\begin{algorithme}[title=\textbf{Branch and Bound}]
\begin{enumerate}
  \item Solve the LP relaxation of the original problem.
  \item If the solution is integral, it is optimal. \textbf{Stop}.
  \item Otherwise, choose a fractional variable $x_j$ (value $\bar{x}_j$).
  \item \textbf{Branching}: create two subproblems:
    \begin{itemize}
      \item Left branch: add $x_j \leq \lfloor \bar{x}_j \rfloor$.
      \item Right branch: add $x_j \geq \lceil \bar{x}_j \rceil$.
    \end{itemize}
  \item \textbf{Bounding}: solve the LP relaxation of each subproblem.
  \item \textbf{Pruning}: eliminate a node if:
    \begin{itemize}
      \item the subproblem is infeasible,
      \item its lower bound $\geq$ best known solution (incumbent),
      \item its solution is integral (update incumbent).
    \end{itemize}
  \item Repeat until all nodes are pruned.
\end{enumerate}
\end{algorithme}

\begin{example}[Branch and Bound]
Consider the ILP:
\[
  \max\; 5x_1 + 4x_2 \quad \text{s.t.} \quad
  6x_1 + 4x_2 \leq 24, \quad x_1 + 2x_2 \leq 6, \quad
  x_1, x_2 \in \Z_+.
\]
The LP relaxation gives $x_1 = 3.6$, $x_2 = 0.6$, $z = 20.4$.
Branching on $x_1$:
\begin{itemize}
  \item Branch $x_1 \leq 3$: LP optimum $x_1 = 3, x_2 = 1.5$, $z = 21$.
  \item Branch $x_1 \geq 4$: LP optimum $x_1 = 4, x_2 = 0$, $z = 20$
        (integral).
\end{itemize}
Continuing on the left branch by branching on $x_2$: $x_2 \leq 1$ gives
$z = 19$ (integral); $x_2 \geq 2$ eventually yields $z = 18$. The optimum
is $x_1 = 4, x_2 = 0$ with $z^* = 20$.
\end{example}

\section{Cutting plane methods}

\begin{definition}[Valid inequality]
\index{cutting plane}
A \textbf{valid inequality} (or cut) is a linear constraint satisfied by all
feasible integer solutions but violated by the current fractional LP solution.
\end{definition}

\subsection{Gomory cuts}

\begin{theorem}[Gomory cut]
\index{Gomory cuts}
Let $x_i = \bar{b}_i - \sum_j \bar{a}_{ij} x_j$ be a row of the final
simplex tableau with $\bar{b}_i \notin \Z$. Then the inequality:
\[
  \sum_j \{\bar{a}_{ij}\} \, x_j \geq \{\bar{b}_i\}
\]
is a valid cut, where $\{y\} = y - \lfloor y \rfloor$ is the fractional part.
\end{theorem}

\begin{proof}
From the simplex row:
$x_i + \sum_j \bar{a}_{ij} x_j = \bar{b}_i$. Decompose into integer and
fractional parts:
\[
  \underbrace{x_i + \sum_j \lfloor \bar{a}_{ij} \rfloor x_j
  - \lfloor \bar{b}_i \rfloor}_{\in \Z}
  = \{\bar{b}_i\} - \sum_j \{\bar{a}_{ij}\} x_j.
\]
The left side is integer. Since $\{\bar{b}_i\} > 0$ and each
$\{\bar{a}_{ij}\} x_j \geq 0$, we obtain
$\sum_j \{\bar{a}_{ij}\} x_j \geq \{\bar{b}_i\}$.
\end{proof}

\begin{remark}
In practice, \textbf{Branch and Cut} methods combine B\&B with cut
generation at each node. Modern solvers (Gurobi, CPLEX) use many families
of cuts beyond Gomory: mixed-integer rounding, lift-and-project, Chv\'atal
cuts, etc.
\end{remark}

\section{Total unimodularity}

\begin{definition}[Totally unimodular matrix]
\index{total unimodularity}
A matrix $A \in \Z^{m \times n}$ is \textbf{totally unimodular} (TU) if every
square submatrix has determinant in $\{-1, 0, 1\}$.
\end{definition}

\begin{theorem}[Exact relaxation]
If $A$ is TU and $b \in \Z^m$, then every vertex of the polyhedron
$\{x \in \R^n : Ax \leq b, x \geq 0\}$ is integral. Hence the LP relaxation
automatically yields an integer solution.
\end{theorem}

\begin{proof}
Every vertex solves a system $A' x = b'$ where $A'$ is an invertible square
submatrix. Since $A$ is TU, $\det(A') \in \{-1, 1\}$, so by Cramer's rule,
$x = (A')^{-1} b'$ is integral.
\end{proof}

\begin{example}[Classical TU matrices]
\begin{itemize}
  \item The vertex-arc incidence matrix of a directed graph is TU.
  \item The constraint matrix of the transportation problem is TU.
  \item Consequence: flow and transportation problems are solved by the simplex
        and automatically yield integer solutions.
\end{itemize}
\end{example}

\section{Classical applications}

\subsection{Knapsack problem}

\begin{definition}[0-1 Knapsack]
\index{knapsack problem}
Given $n$ items with values $v_1, \ldots, v_n$ and weights $w_1, \ldots, w_n$,
and a knapsack of capacity $W$:
\[
  \max \sum_{j=1}^n v_j x_j \quad \text{s.t.} \quad
  \sum_{j=1}^n w_j x_j \leq W, \quad x_j \in \{0, 1\}.
\]
\end{definition}

\subsection{Traveling salesman problem}

\begin{definition}[TSP formulation]
\index{traveling salesman problem}
The \textbf{Traveling Salesman Problem} (TSP) on $n$ cities:
\[
  \min \sum_{i \ne j} d_{ij} x_{ij} \quad \text{s.t.} \quad
  \sum_{j \ne i} x_{ij} = 1 \;\forall i, \quad
  \sum_{i \ne j} x_{ij} = 1 \;\forall j, \quad x_{ij} \in \{0,1\},
\]
plus subtour elimination constraints (e.g., Miller-Tucker-Zemlin or
exponentially many subtour cuts).
\end{definition}

\begin{warning}[title=\textbf{TSP complexity}]
TSP is NP-hard. State-of-the-art solvers (Concorde) use advanced Branch and
Cut and can solve instances with several thousand cities to optimality.
\end{warning}

\section{Exercises}

\begin{exercise}[$\star$ -- Branch and Bound]
Solve by B\&B: $\max\; 3x_1 + 2x_2$ s.t. $2x_1 + x_2 \leq 6$,
$x_1 + 3x_2 \leq 9$, $x_1, x_2 \in \Z_+$. Draw the search tree.
\end{exercise}

\begin{exercise}[$\star\star$ -- Gomory cuts]
Consider $\max\; x_1 + x_2$ s.t. $3x_1 + 2x_2 \leq 6$,
$x_1 + 4x_2 \leq 4$, $x_1, x_2 \geq 0$ integer. Generate a Gomory cut from
the final simplex tableau and solve.
\end{exercise}

\begin{exercise}[$\star\star$ -- Knapsack]
Solve the 0-1 knapsack with values $(10, 6, 12, 7)$, weights $(3, 2, 5, 3)$,
capacity $W = 8$ by Branch and Bound.
\end{exercise}

\begin{exercise}[$\star\star\star$ -- TSP formulation]
Formulate TSP on 5 cities with MTZ constraints. Solve the LP relaxation and
measure the integrality gap.
\end{exercise}

\begin{exercise}[$\star\star\star$ -- Total unimodularity]
Prove that the vertex-arc incidence matrix of a directed graph is TU. Deduce
that the minimum cost flow problem always has an integer optimal solution when
the data are integral.
\end{exercise}
